\begin{lemma}
Under the hypotheses of \(\noderef{HugeFamilySizeBound}\), and assuming also
\(\varepsilon_1\ge0\), suppose the sample parameters are the rounded hierarchy
parameters
\[
m=\noderef{TwoBiteNaturalReducedVertexCount}(n),
\qquad
p=\noderef{TwoBiteEdgeProbability}\!\left(\frac12,n\right).
\]
Then the opposite-color huge projected-pair sums satisfy
\[
\sum_{x\in H_I\cap V_R}\binom{|\pi_B(X_x(I))|}{2}_{\mathbb R}
\le
(1+\varepsilon_1)
\min\left\{
\binom{k-|\pi_R(I)|}{2}_{\mathbb R},
\binom{pn}{2}_{\mathbb R}+
\binom{k-|\pi_R(I)|-pn}{2}_{\mathbb R}
\right\},
\]
and the color-swapped estimate with \(R\) and \(B\) interchanged.
\end{lemma}
\begin{proof}
Let
\[
C(u)=\noderef{RealChooseTwo}(u)=\binom{u}{2}_{\mathbb R}.
\]
We first record the two finite convexity estimates used below.
In each application the finite family is indexed by a Lean `Finset.filter`;
choosing any enumeration of that finite set turns the sum into the list form
written here, and the estimates are invariant under the chosen enumeration.

First, if \(b_1,\ldots,b_h\) are nonnegative real numbers, then
\[
\sum_i C(b_i)\le C\left(\sum_i b_i\right). \tag{1}
\]
For two nonnegative terms this is the identity
\[
C(r+s)-C(r)-C(s)=rs\ge0.
\]
Induct on the finite list: the partial sum already accumulated is
nonnegative, the next term is nonnegative, and the two-term identity shows
that replacing \(C(P)+C(b)\) by \(C(P+b)\) can only increase the right side.
This proves (1), including the empty-list case where both sides are \(0\).

Second, suppose \(0\le M\), \(0\le b_i\le M\) for every \(i\), let
\[
B=\sum_i b_i,
\]
and assume \(M\le B\).  Then
\[
\sum_i C(b_i)\le C(M)+C(B-M). \tag{2}
\]
If the list is empty, then \(0\le M\le B=0\), so \(M=0\) and both sides are
\(0\).  Otherwise
order the finite list arbitrarily and let \(j\) be the first index for which
the partial sum \(P_j=\sum_{i\le j}b_i\) is at least \(M\).  Then
\[
P_{j-1}\le M,\qquad b_j\le M,\qquad P_j-M\ge0.
\]
By (1) applied to the first \(j-1\) terms,
\[
\sum_{i<j}C(b_i)\le C(P_{j-1}).
\]
Since \(P_j=P_{j-1}+b_j\), the explicit difference formula
\[
C(M)+C(P_j-M)-C(P_{j-1})-C(b_j)
  =(M-P_{j-1})(M-b_j)
\]
is nonnegative.  Thus
\[
\sum_{i\le j}C(b_i)\le C(M)+C(P_j-M).
\]
All remaining \(b_i\) and \(P_j-M\) are nonnegative, so applying (1) to
\(P_j-M\) together with the remaining terms gives
\[
C(P_j-M)+\sum_{i>j}C(b_i)
\le C(B-M).
\]
Adding this to the preceding display proves (2).  In the complementary case
\(B\le M\), (1) gives directly
\[
\sum_i C(b_i)\le C(B). \tag{3}
\]

We now prove the red-to-blue estimate and match it to the first Lean
conjunct.  Define the exact huge-red predicate and filter by
\[
\mathsf{hugeRed}(x)\iff
\noderef{TwoBiteHugeClass}(\omega,I,x)\wedge
  \exists r,\ x=\operatorname{inl}(r),
\]
\[
H_R=\operatorname{univ}.\operatorname{filter}(\mathsf{hugeRed}).
\]
For \(x\in H_R\), set
\[
a_x=
\bigl|X_x(I).\operatorname{image}
  (\noderef{TwoBiteBlueProjection}(\omega))\bigr|,
\]
and put
\[
T_R=\sum_{x\in H_R}a_x,\qquad
D_R=(k:\mathbb R)-
  \bigl((I.\operatorname{image}
    (\noderef{TwoBiteRedProjection}(\omega))).\operatorname{card}:\mathbb R\bigr),
\]
\[
b=p\cdot(n:\mathbb R),\qquad M=(1+\varepsilon_2)b.
\]
Unfolding \(\noderef{TwoBiteProjectedPairSum}\), the left side of the first
Lean conjunct is exactly
\[
S_R=\sum_{x\in H_R}C(a_x)
=\noderef{TwoBiteProjectedPairSum}
  (\omega,I,\mathsf{hugeRed},\noderef{TwoBiteBlueProjection}(\omega)). \tag{4}
\]
Notice the deliberate cross-color bookkeeping: the summand uses the blue
projection \(\pi_B\), but the deficit \(D_R\) uses
\(|I.\operatorname{image}\pi_R|\), exactly as in the Lean statement.

Apply \(\noderef{HugeOppositeProjectionMassControl}\) with the same
\(\omega,I,n,m,k,p,\varepsilon,\varepsilon_1,\varepsilon_2,n_0\), the same
hypotheses \(0\le\varepsilon_1\),
\(\noderef{ParameterHierarchyEventualComparisons}\), \(n_0<n\), the same
rounded equalities for \(m,p,k\), the same \(|I|=k\), and the same
\(\noderef{FiberAndDegreeEvent}(\omega,\varepsilon_2)\).  In the notation of
that child, its `redMass' is exactly \(T_R\), its `redDeficit' is exactly
\(D_R\), and its `blockScale' is exactly \(b=p(n:\mathbb R)\).  Its
first, third, fourth, and fifth red conjuncts give respectively
\[
0\le a_x\le M\qquad(x\in H_R), \tag{5}
\]
where nonnegativity is cardinality nonnegativity, and
\[
C(T_R)\le(1+\varepsilon_1)C(D_R), \tag{6}
\]
\[
T_R\le M\Longrightarrow
C(T_R)\le
(1+\varepsilon_1)\bigl(C(b)+C(D_R-b)\bigr), \tag{7}
\]
\[
M\le T_R\Longrightarrow
C(M)+C(T_R-M)\le
(1+\varepsilon_1)\bigl(C(b)+C(D_R-b)\bigr). \tag{8}
\]
The unused second red conjunct is the mass-deficit estimate
\(T_R\le(1+\varepsilon_2)D_R\); it is part of the child package but is not
needed for this aggregation step.
The hierarchy package gives \(0\le\varepsilon_2\), and the rounded definition
of \(p\) gives \(0\le b\), so \(0\le M\).  Hence the convexity estimates above
apply to the finite family \((a_x)_{x\in H_R}\).

By (1) and (6),
\[
S_R=\sum_{x\in H_R}C(a_x)
\le C(T_R)
\le(1+\varepsilon_1)C(D_R). \tag{9}
\]
This is the one-block side of the minimum.  For the two-block side, split the
linear order on \(T_R\) and \(M\).  If \(T_R\le M\), then (3) and (7) give
\[
S_R\le C(T_R)
\le(1+\varepsilon_1)\bigl(C(b)+C(D_R-b)\bigr). \tag{10}
\]
If \(M\le T_R\), then (2), the pointwise cap (5), and (8) give
\[
S_R\le C(M)+C(T_R-M)
\le(1+\varepsilon_1)\bigl(C(b)+C(D_R-b)\bigr). \tag{11}
\]
The two cases cover all possibilities, so (10)--(11) prove the two-block side
unconditionally:
\[
S_R\le(1+\varepsilon_1)\bigl(C(b)+C(D_R-b)\bigr). \tag{12}
\]
Let \(A_R=C(D_R)\) and \(B_R=C(b)+C(D_R-b)\).  From (9) and (12) we have
\[
S_R\le(1+\varepsilon_1)A_R,\qquad
S_R\le(1+\varepsilon_1)B_R.
\]
By the defining property of the minimum in a linear order, these two
inequalities give
\[
S_R\le\min\{(1+\varepsilon_1)A_R,(1+\varepsilon_1)B_R\}.
\]
Because \(1+\varepsilon_1\ge0\), multiplication by \(1+\varepsilon_1\) is
monotone and
\[
\min\{(1+\varepsilon_1)A_R,(1+\varepsilon_1)B_R\}
=(1+\varepsilon_1)\min\{A_R,B_R\}.
\]
Combining these two displays yields
\[
S_R\le
(1+\varepsilon_1)\min\{C(D_R),\,C(b)+C(D_R-b)\},
\]
which is precisely the red-huge/blue-projection inequality in the Lean
statement, with \(b=p\cdot(n:\mathbb R)\).

We prove the blue-to-red estimate separately, again matching the Lean
conjunct.  Define
\[
\mathsf{hugeBlue}(x)\iff
\noderef{TwoBiteHugeClass}(\omega,I,x)\wedge
  \exists q,\ x=\operatorname{inr}(q),
\]
\[
H_B=\operatorname{univ}.\operatorname{filter}(\mathsf{hugeBlue}).
\]
For \(x\in H_B\), set
\[
c_x=
\bigl|X_x(I).\operatorname{image}
  (\noderef{TwoBiteRedProjection}(\omega))\bigr|,
\]
and put
\[
T_B=\sum_{x\in H_B}c_x,\qquad
D_B=(k:\mathbb R)-
  \bigl((I.\operatorname{image}
    (\noderef{TwoBiteBlueProjection}(\omega))).\operatorname{card}:\mathbb R\bigr).
\]
Unfolding \(\noderef{TwoBiteProjectedPairSum}\), the left side of the second
Lean conjunct is
\[
S_B=\sum_{x\in H_B}C(c_x)
=\noderef{TwoBiteProjectedPairSum}
  (\omega,I,\mathsf{hugeBlue},\noderef{TwoBiteRedProjection}(\omega)). \tag{13}
\]
Here the summand uses the red projection \(\pi_R\), while the deficit \(D_B\)
uses \(|I.\operatorname{image}\pi_B|\), exactly as required.

The blue conclusions of
\(\noderef{HugeOppositeProjectionMassControl}\), instantiated with the same
data as above, are read with `blueMass' \(=T_B\), `blueDeficit' \(=D_B\), and
the same `blockScale' \(=b\).  Its first, third, fourth, and fifth blue
conjuncts give respectively
\[
0\le c_x\le M\qquad(x\in H_B), \tag{14}
\]
\[
C(T_B)\le(1+\varepsilon_1)C(D_B), \tag{15}
\]
\[
T_B\le M\Longrightarrow
C(T_B)\le
(1+\varepsilon_1)\bigl(C(b)+C(D_B-b)\bigr), \tag{16}
\]
and
\[
M\le T_B\Longrightarrow
C(M)+C(T_B-M)\le
(1+\varepsilon_1)\bigl(C(b)+C(D_B-b)\bigr). \tag{17}
\]
Again, the child also supplies \(T_B\le(1+\varepsilon_2)D_B\), but the
present proof only needs the displayed pointwise and binomial conclusions.
Using (1) and (15),
\[
S_B\le C(T_B)\le(1+\varepsilon_1)C(D_B). \tag{18}
\]
For the two-block side, if \(T_B\le M\), then (3) and (16) give
\[
S_B\le(1+\varepsilon_1)\bigl(C(b)+C(D_B-b)\bigr),
\]
while if \(M\le T_B\), then (2), the pointwise cap (14), and (17) give the
same inequality.  Thus
\[
S_B\le(1+\varepsilon_1)\bigl(C(b)+C(D_B-b)\bigr). \tag{19}
\]
Let \(A_B=C(D_B)\) and \(B_B=C(b)+C(D_B-b)\).  Equations (18) and (19) give
\[
S_B\le(1+\varepsilon_1)A_B,\qquad
S_B\le(1+\varepsilon_1)B_B.
\]
Hence
\[
S_B\le\min\{(1+\varepsilon_1)A_B,(1+\varepsilon_1)B_B\}.
\]
As above, \(1+\varepsilon_1\ge0\), so
\[
\min\{(1+\varepsilon_1)A_B,(1+\varepsilon_1)B_B\}
=(1+\varepsilon_1)\min\{A_B,B_B\}.
\]
Therefore
\[
S_B\le
(1+\varepsilon_1)\min\{C(D_B),\,C(b)+C(D_B-b)\}.
\]
This is exactly the blue-huge/red-projection inequality in the second Lean
conjunct.  The two proved inequalities are the desired conjunction.
\end{proof}
